%=====================================================================
%  Simulando Qubits com PySCF: Uma Jornada do Zero aos Pontos
%  Quanticos para Computacao Quantica
%  Autores: Cristiano da Costa Alves
%           Nilseia Aparecida Barbosa
%           Fernando Manuel Araujo Moreira
%
%  Arquivo principal do livro. Compile com:
%      latexmk -pdf Livro_PySCF.tex
%  ou  pdflatex Livro_PySCF.tex  (2x, para sumario/referencias)
%=====================================================================
\documentclass[11pt,a4paper,openany]{book}

%--------------------- Codificacao e idioma --------------------------
\usepackage[utf8]{inputenc}
\usepackage[T1]{fontenc}
\usepackage[brazil]{babel}
\usepackage{lmodern}
\usepackage[zerostyle=b,scaled=.85]{newtxtt} % monospace legivel
\usepackage{microtype}

%--------------------- Matematica e graficos -------------------------
\usepackage{amsmath,amssymb,amsfonts}
\usepackage{mathtools}
\usepackage{graphicx}
\usepackage{booktabs}
\usepackage{multirow}
\usepackage{array}
\usepackage{enumitem}

%--------------------- Layout de pagina ------------------------------
\usepackage{geometry}
\geometry{left=3cm,right=2.5cm,top=2.8cm,bottom=2.8cm}
\usepackage{emptypage}
\linespread{1.05}

%--------------------- Cabecalhos ------------------------------------
\usepackage{fancyhdr}
\pagestyle{fancy}
\fancyhf{}
\fancyhead[LE]{\small\itshape\nouppercase{\leftmark}}
\fancyhead[RO]{\small\itshape\nouppercase{\rightmark}}
\fancyfoot[C]{\thepage}
\renewcommand{\headrulewidth}{0.4pt}

%--------------------- Cores -----------------------------------------
\usepackage{xcolor}
\definecolor{azulpyscf}{RGB}{18,72,120}
\definecolor{azulclaro}{RGB}{232,240,248}
\definecolor{cinzacodigo}{RGB}{247,247,247}
\definecolor{verdecom}{RGB}{60,120,60}
\definecolor{roxokw}{RGB}{140,40,140}
\definecolor{laranjastr}{RGB}{170,80,20}
\definecolor{verdeok}{RGB}{20,110,60}
\definecolor{vermelhoaviso}{RGB}{150,30,30}

%--------------------- Capitulos estilizados -------------------------
\usepackage[Sonny]{fncychap}
\ChNameVar{\Large\sffamily\bfseries\color{azulpyscf}}
\ChNumVar{\Huge\sffamily\bfseries\color{azulpyscf}}
\ChTitleVar{\Large\sffamily\bfseries\color{azulpyscf}}

%--------------------- Codigo Python (listings) ----------------------
\usepackage{listings}
\lstset{
  basicstyle=\ttfamily\small,
  keywordstyle=\color{roxokw}\bfseries,
  commentstyle=\color{verdecom}\itshape,
  stringstyle=\color{laranjastr},
  numberstyle=\tiny\color{gray},
  numbers=left,
  numbersep=8pt,
  backgroundcolor=\color{cinzacodigo},
  frame=single,
  rulecolor=\color{gray!40},
  framesep=6pt,
  breaklines=true,
  showstringspaces=false,
  tabsize=4,
  columns=fullflexible,
  keepspaces=true,
  inputencoding=utf8,
  extendedchars=true,
  literate=%
    {á}{{\'a}}1 {é}{{\'e}}1 {í}{{\'i}}1 {ó}{{\'o}}1 {ú}{{\'u}}1
    {Á}{{\'A}}1 {É}{{\'E}}1 {Í}{{\'I}}1 {Ó}{{\'O}}1 {Ú}{{\'U}}1
    {â}{{\^a}}1 {ê}{{\^e}}1 {ô}{{\^o}}1 {Â}{{\^A}}1 {Ê}{{\^E}}1
    {ã}{{\~a}}1 {õ}{{\~o}}1 {Ã}{{\~A}}1 {Õ}{{\~O}}1
    {à}{{\`a}}1 {ç}{{\c c}}1 {Ç}{{\c C}}1 {º}{{\textordmasculine}}1
    {ª}{{\textordfeminine}}1
}
\lstdefinestyle{python}{language=Python}
\lstdefinestyle{console}{
  language=bash,
  numbers=none,
  keywordstyle=\color{black},
  backgroundcolor=\color{black!88},
  basicstyle=\ttfamily\small\color{white},
  frame=single,
  rulecolor=\color{black},
}

%--------------------- Caixas didaticas ------------------------------
\usepackage[most]{tcolorbox}

\newtcolorbox{caixanota}{
  colback=azulclaro, colframe=azulpyscf, coltitle=white,
  fonttitle=\bfseries\sffamily, title=Nota,
  boxrule=0.6pt, arc=2pt, left=6pt,right=6pt,top=4pt,bottom=4pt,
  breakable, enhanced}

\newtcolorbox{caixadica}{
  colback=green!5, colframe=verdeok, coltitle=white,
  fonttitle=\bfseries\sffamily, title=Dica,
  boxrule=0.6pt, arc=2pt, left=6pt,right=6pt,top=4pt,bottom=4pt,
  breakable, enhanced}

\newtcolorbox{caixaatencao}{
  colback=red!5, colframe=vermelhoaviso, coltitle=white,
  fonttitle=\bfseries\sffamily, title=Atenção,
  boxrule=0.6pt, arc=2pt, left=6pt,right=6pt,top=4pt,bottom=4pt,
  breakable, enhanced}

\newtcolorbox{caixamaos}[1][Mãos à obra]{
  colback=orange!6, colframe=orange!70!black, coltitle=white,
  fonttitle=\bfseries\sffamily, title={#1},
  boxrule=0.8pt, arc=2pt, left=6pt,right=6pt,top=4pt,bottom=4pt,
  breakable, enhanced}

\newtcolorbox{caixaancora}{
  colback=azulpyscf!8, colframe=azulpyscf, coltitle=white,
  fonttitle=\bfseries\sffamily, title={Projeto Âncora},
  boxrule=1pt, arc=2pt, left=6pt,right=6pt,top=4pt,bottom=4pt,
  breakable, enhanced}

%--------------------- Comandos auxiliares ---------------------------
\newcommand{\code}[1]{\texttt{#1}}
\newcommand{\pyscf}{\textsf{PySCF}}
\newcommand{\numpy}{\textsf{NumPy}}
\newcommand{\python}{\textsf{Python}}
\newcommand{\qubit}{\emph{qubit}}
\newcommand{\Mole}{\texttt{Mole}}

%--------------------- Hyperref (por ultimo) -------------------------
\usepackage[hidelinks,colorlinks=true,linkcolor=azulpyscf,
            urlcolor=azulpyscf,citecolor=azulpyscf]{hyperref}
